\documentclass[11pt,a4paper]{article}
\usepackage[utf8]{inputenc}
\usepackage[margin=1in]{geometry}
\usepackage{amsmath}
\usepackage{amssymb}
\usepackage{enumitem}
\usepackage{booktabs}
\usepackage{titlesec}
\usepackage{microtype}
\usepackage{xcolor}

\titleformat{\section}{\large\bfseries\sffamily}{}{0em}{}[\titlerule]
\titleformat{\subsection}{\normalsize\bfseries\sffamily}{}{0em}{}

\setlength{\parindent}{0pt}
\setlength{\parskip}{8pt plus 2pt minus 1pt}

\title{\textbf{CAMS Exam: Expanded Section 2 — International Standards}}
\author{\textbf{Based on Candidate Feedback (FATF \& PATRIOT Act Heaviest)}}
\date{\today}

\begin{document}
	
	\maketitle
	
	\section*{IMPORTANT NOTE FROM CANDIDATES}
	Based on recent exam-taker feedback:
	\begin{itemize}
		\item \textbf{FATF is the single most heavily tested topic} — know Recommendations, Mutual Evaluations, Immediate Outcomes, and Grey/Black Lists
		\item \textbf{USA PATRIOT Act appears repeatedly} — Sections 311, 313, 314(a), 314(b), 319(b) are essential
		\item \textbf{Most questions require selecting 2-4 answers} — you must get ALL correct to get credit
		\item \textbf{Global regulations section is the most difficult} — focus your study time here
	\end{itemize}
	
	\section*{Instructions}
	Each question has \textbf{five options (A through E)}. 
	The question will specify how many answers are correct: \textbf{[SELECT TWO]}, \textbf{[SELECT THREE]}, or \textbf{[SELECT FOUR]}.
	
	\newpage
	
	% ============================================================
	\section{FATF Recommendations (Deep Coverage)}
	% ============================================================
	
	\subsection{1. FATF Recommendation 1: Risk-Based Approach \textnormal{[SELECT TWO]}}
	
	A national AML authority is developing its AML/CFT framework. It must decide whether to apply a risk-based approach (RBA) or a prescriptive rules-based approach.
	
	Which two statements correctly describe the RBA as required by FATF Recommendation 1?
	
	\begin{enumerate}[label=\Alph*.]
		\item The RBA requires countries and financial institutions to identify, assess, and understand their ML/TF risks and apply measures commensurate with those risks.
		\item The RBA allows simplified measures where risks are lower, provided that the simplified measures are commensurate with the nature of lower risk.
		\item The RBA requires the same level of due diligence for all customers regardless of risk profile to ensure consistent treatment.
		\item The RBA eliminates all due diligence requirements for customers from FATF member countries.
		\item Under the RBA, enhanced measures are always required for all cross-border correspondent banking relationships without exception.
	\end{enumerate}
	
	\textbf{Correct Answers: A and B}
	
	\subsection{2. FATF Recommendation 2: National Cooperation and Coordination \textnormal{[SELECT TWO]}}
	
	A FATF assessor is reviewing a country's implementation of Recommendation 2. The country has multiple agencies responsible for AML/CFT: a financial intelligence unit, a financial supervisor, a police financial crimes unit, and a prosecutor's office.
	
	Which two findings would indicate non-compliance with Recommendation 2?
	
	\begin{enumerate}[label=\Alph*.]
		\item The FIU receives STRs but has never disseminated financial intelligence to law enforcement for use in an investigation.
		\item The supervisor has identified significant AML deficiencies at a bank but has not communicated these findings to the FIU or law enforcement.
		\item The country has established a formal inter-agency coordination committee that meets quarterly.
		\item The police and prosecutors have a formal memorandum of understanding for sharing evidence.
		\item The central bank publishes annual AML/CFT supervision reports on its website.
	\end{enumerate}
	
	\textbf{Correct Answers: A and B}
	
	\subsection{3. FATF Recommendation 3: Money Laundering Offense \textnormal{[SELECT TWO]}}
	
	A country is drafting its money laundering legislation and must determine the scope of the offense.
	
	Which two statements correctly describe FATF's requirements for the money laundering offense under Recommendation 3?
	
	\begin{enumerate}[label=\Alph*.]
		\item The offense must apply to all serious predicate offenses, with FATF recommending the "all-crimes approach."
		\item The offense must apply to persons who committed the predicate offense (self-laundering is criminalized).
		\item The offense only applies if the predicate offense was committed in a different country.
		\item The offense must require proof that the launderer knew the exact predicate offense that generated the proceeds.
		\item The offense only applies to drug trafficking and terrorism financing proceeds.
	\end{enumerate}
	
	\textbf{Correct Answers: A and B}
	
	\subsection{4. FATF Recommendation 5: Terrorist Financing Offense \textnormal{[SELECT TWO]}}
	
	A country has criminalized terrorist financing only when the funds are used for a specific terrorist act, and only when the amount exceeds \$10,000. The country has no TF convictions in ten years.
	
	Which two statements correctly describe the deficiencies under Recommendation 5?
	
	\begin{enumerate}[label=\Alph*.]
		\item TF must be criminalized regardless of whether the funds are actually used to carry out a terrorist act; collection of funds for terrorist purposes is sufficient.
		\item TF should be criminalized without any minimum monetary threshold; a \$10,000 threshold violates Recommendation 5.
		\item The \$10,000 threshold is acceptable as long as the country files STRs for TF below that amount.
		\item TF does not need to be criminalized as a standalone offense if it is covered under money laundering statutes.
		\item The zero convictions are irrelevant; only the existence of the law matters for technical compliance.
	\end{enumerate}
	
	\textbf{Correct Answers: A and B}
	
	\subsection{5. FATF Recommendation 8: Non-Profit Organizations \textnormal{[SELECT TWO]}}
	
	A country has enacted no specific AML/CFT measures for the non-profit organization sector. A FATF assessor notes that the country has 500 NPOs, 40 of which operate in conflict zones where terrorist groups are active, and 15 of which have had their bank accounts closed due to suspicious activity.
	
	Which two statements correctly describe the country's obligations under Recommendation 8?
	
	\begin{enumerate}[label=\Alph*.]
		\item The country must apply targeted risk-based supervision of NPOs that are vulnerable to terrorist financing abuse.
		\item The country must have mechanisms to ensure that NPOs are not misused by terrorist organizations, including through outreach, supervision, and monitoring of cross-border transactions.
		\item All NPOs must register with the FIU regardless of their size or risk profile.
		\item NPOs operating in conflict zones are exempt from all AML/CFT requirements because of humanitarian access concerns.
		\item Recommendation 8 only applies to NPOs that receive government funding.
	\end{enumerate}
	
	\textbf{Correct Answers: A and B}
	
	\subsection{6. FATF Recommendation 10: Customer Due Diligence \textnormal{[SELECT THREE]}}
	
	A financial institution is implementing its CDD program. The compliance officer must ensure all required elements are covered.
	
	Which three of the following are required CDD measures under Recommendation 10?
	
	\begin{enumerate}[label=\Alph*.]
		\item Identifying the customer and verifying their identity using reliable, independent source documents or information.
		\item Identifying the beneficial owner and taking reasonable measures to verify their identity.
		\item Understanding the purpose and intended nature of the business relationship.
		\item Obtaining a credit report for every customer regardless of risk level.
		\item Conducting ongoing due diligence on the business relationship and scrutinizing transactions throughout the relationship.
	\end{enumerate}
	
	\textbf{Correct Answers: A, B, and C}
	
	\subsection{7. FATF Recommendation 11: Recordkeeping \textnormal{[SELECT ONE]}}
	
	A bank's compliance officer is reviewing the recordkeeping policy to ensure compliance with FATF Recommendation 11.
	
	What is the minimum required retention period for records under Recommendation 11?
	
	\begin{enumerate}[label=\Alph*.]
		\item 3 years from the date of transaction or account closure.
		\item 5 years from the date of transaction or account closure.
		\item 7 years from the date of transaction or account closure.
		\item 10 years from the date of transaction or account closure.
		\item Indefinitely for all records related to SAR filings.
	\end{enumerate}
	
	\textbf{Correct Answer: B}
	
	\subsection{8. FATF Recommendation 12: DNFBPs — Casinos \textnormal{[SELECT TWO]}}
	
	A country is implementing CDD requirements for designated non-financial businesses and professions (DNFBPs). The country must determine which obligations apply to casinos.
	
	Which two CDD obligations apply to casinos under Recommendation 12?
	
	\begin{enumerate}[label=\Alph*.]
		\item Casinos must conduct CDD when a customer purchases chips or tokens equivalent to USD/EUR 3,000 or more.
		\item Casinos must conduct CDD for any transaction exceeding USD/EUR 10,000.
		\item Casinos must verify the identity of customers and the source of funds for large transactions.
		\item Casinos are exempt from CDD requirements because gambling is regulated separately.
		\item Casinos must file CTRs but not SARs under Recommendation 12.
	\end{enumerate}
	
	\textbf{Correct Answers: A and C}
	
	\subsection{9. FATF Recommendation 14: Money or Value Transfer Services \textnormal{[SELECT TWO]}}
	
	A country has many unlicensed money remitters (hawala operators) operating informally. The government has decided not to regulate them because they serve migrant workers who cannot access formal banking.
	
	Which two statements correctly describe the country's obligations under Recommendation 14?
	
	\begin{enumerate}[label=\Alph*.]
		\item All persons that provide money or value transfer services must be licensed or registered and subject to AML/CFT monitoring.
		\item Unlicensed MVTS providers are subject to criminal, civil, or administrative sanctions.
		\item Hawala operators are exempt from licensing requirements because of cultural heritage protections under FATF.
		\item Remittances below \$1,000 are exempt from CDD requirements under Recommendation 14.
		\item Countries may exempt MVTS providers if they process less than \$1 million annually.
	\end{enumerate}
	
	\textbf{Correct Answers: A and B}
	
	\subsection{10. FATF Recommendation 15: New Technologies \textnormal{[SELECT TWO]}}
	
	A country has a growing virtual asset sector but has not yet implemented AML/CFT requirements for VASPs. The FATF assessor notes that virtual assets are being used in ML/TF cases.
	
	Which two statements correctly describe the country's obligations under Recommendation 15?
	
	\begin{enumerate}[label=\Alph*.]
		\item Countries must assess the ML/TF risks associated with new technologies and virtual assets.
		\item VASPs must be licensed or registered and subject to AML/CFT obligations.
		\item Virtual assets are exempt from the Travel Rule because of technical limitations.
		\item Only fiat-to-crypto exchanges must be regulated; crypto-to-crypto exchanges are exempt.
		\item DeFi protocols are automatically exempt from regulation regardless of control structure.
	\end{enumerate}
	
	\textbf{Correct Answers: A and B}
	
	\subsection{11. FATF Recommendation 16: Wire Transfer Rule (Travel Rule) \textnormal{[SELECT THREE]}}
	
	An international payments processor is auditing its compliance with the Wire Transfer Rule. The audit covers cross-border wire transfers.
	
	Which three pieces of information must accompany a cross-border wire transfer of USD/EUR 2,000?
	
	\begin{enumerate}[label=\Alph*.]
		\item Originator's full name.
		\item Originator's account number (or unique transaction reference number).
		\item Originator's address, national ID number, customer ID number, or date and place of birth (at least one of these).
		\item Originator's favorite color or hobby.
		\item Originator's mother's maiden name.
	\end{enumerate}
	
	\textbf{Correct Answers: A, B, and C}
	
	\subsection{12. FATF Recommendation 19: Higher-Risk Jurisdictions \textnormal{[SELECT TWO]}}
	
	A FATF member country is considering its response to a jurisdiction designated as high-risk (Black List) by FATF.
	
	Which two statements correctly describe the country's obligations under Recommendation 19?
	
	\begin{enumerate}[label=\Alph*.]
		\item Countries must apply enhanced due diligence measures to business relationships and transactions with natural and legal persons from high-risk jurisdictions.
		\item FATF calls on members to apply countermeasures, which may include prohibiting correspondent relationships, restricting transactions, or requiring enhanced reporting.
		\item Countermeasures are optional and can be ignored if the member country disagrees with FATF's designation.
		\item Black List designations are purely advisory and have no legal effect on financial institutions.
		\item Only the United Nations Security Council can designate high-risk jurisdictions; FATF designations are non-binding.
	\end{enumerate}
	
	\textbf{Correct Answers: A and B}
	
	\subsection{13. FATF Recommendation 20: STR Filing \textnormal{[SELECT ONE]}}
	
	A compliance officer at a bank has identified a transaction that gives rise to reasonable suspicion of money laundering. The transaction is valued at \$2,500.
	
	Under Recommendation 20, what is the bank's obligation?
	
	\begin{enumerate}[label=\Alph*.]
		\item The bank must file an STR only if the transaction exceeds \$10,000.
		\item The bank must file an STR immediately upon suspicion, regardless of the transaction value.
		\item The bank must investigate further to obtain conclusive proof before filing an STR.
		\item The bank must notify the customer and give them 30 days to explain the transaction before filing.
		\item The bank must wait 90 days to confirm the suspicious pattern before filing.
	\end{enumerate}
	
	\textbf{Correct Answer: B}
	
	\subsection{14. FATF Recommendation 22: DNFBPs — Lawyers \textnormal{[SELECT TWO]}}
	
	A country is implementing CDD requirements for lawyers and notaries. The legal profession association argues that attorney-client privilege exempts lawyers from all AML obligations.
	
	Which two statements correctly describe the obligations of lawyers under Recommendation 22?
	
	\begin{enumerate}[label=\Alph*.]
		\item Lawyers must conduct CDD when they prepare for or carry out transactions for a client concerning: buying/selling real estate, managing client money/securities, managing bank/savings accounts, or creating/operating legal persons or trusts.
		\item Lawyers are not required to apply CDD when they are providing legal advice about the legal position of their client or representing a client in legal proceedings (including advice on instituting or avoiding legal proceedings).
		\item All lawyer-client communications are fully exempt from AML requirements regardless of the transaction type.
		\item Lawyers never have any CDD obligations under FATF.
		\item Lawyers must file STRs on all real estate transactions regardless of suspicion.
	\end{enumerate}
	
	\textbf{Correct Answers: A and B}
	
	\subsection{15. FATF Recommendation 23: DNFBPs — Trust and Company Service Providers \textnormal{[SELECT TWO]}}
	
	A trust and company service provider (TCSP) is being onboarded by a bank. The TCSP creates shell companies for clients and serves as a nominee director.
	
	Which two obligations apply to TCSPs under Recommendation 23?
	
	\begin{enumerate}[label=\Alph*.]
		\item TCSPs must conduct CDD when they form legal persons or arrange for another person to act as a director or nominee shareholder.
		\item TCSPs must identify beneficial owners of the legal persons they create and maintain records for at least five years.
		\item TCSPs are completely exempt from CDD because they are service providers, not financial institutions.
		\item TCSPs only have CDD obligations if they also provide accounting services.
		\item TCSPs must file SARs but are exempt from CDD under Recommendation 23.
	\end{enumerate}
	
	\textbf{Correct Answers: A and B}
	
	\subsection{16. FATF Recommendation 24: Transparency of Legal Persons \textnormal{[SELECT TWO]}}
	
	A country allows corporations to issue bearer shares without any immobilization or recordkeeping. A FATF assessor identifies this as a significant vulnerability.
	
	Which two statements correctly describe the country's obligations under Recommendation 24?
	
	\begin{enumerate}[label=\Alph*.]
		\item Countries should prohibit bearer shares or require immobilization with a regulated custodian.
		\item Countries must ensure that adequate, accurate, and current information on beneficial ownership of legal persons is obtained and held.
		\item Bearer shares are fully compliant with FATF standards as long as they are physically held by the beneficial owner.
		\item Countries are not required to collect beneficial ownership information; only financial institutions have that obligation.
		\item Bearer shares present no money laundering risk because they are a legitimate corporate governance tool.
	\end{enumerate}
	
	\textbf{Correct Answers: A and B}
	
	\subsection{17. FATF Recommendation 25: Transparency of Legal Arrangements (Trusts) \textnormal{[SELECT TWO]}}
	
	A country has many discretionary trusts but no requirement to identify or register beneficial owners of trusts. A FATF assessor notes that trusts are being used to obscure ownership of real estate.
	
	Which two statements correctly describe the country's obligations under Recommendation 25?
	
	\begin{enumerate}[label=\Alph*.]
		\item Countries must ensure that trustees of express trusts hold adequate, accurate, and current information on beneficial ownership of the trust.
		\item Countries must ensure that trustees disclose their status to financial institutions when forming a business relationship.
		\item Discretionary trusts are exempt from beneficial ownership requirements under Recommendation 25.
		\item Only trusts with assets exceeding \$1 million must disclose beneficial ownership information.
		\item Trusts established in other countries are exempt from all requirements in the country where they operate.
	\end{enumerate}
	
	\textbf{Correct Answers: A and B}
	
	\subsection{18. FATF Recommendation 26: Regulation and Supervision of FIs \textnormal{[SELECT TWO]}}
	
	A country has a financial supervisor that conducts on-site examinations every five years for all banks. The supervisor has not conducted any off-site monitoring or thematic reviews. The supervisor has issued no penalties for AML violations in ten years.
	
	Which two findings would indicate deficiencies under Recommendation 26?
	
	\begin{enumerate}[label=\Alph*.]
		\item Supervisors must have adequate powers to supervise and ensure compliance, including the power to conduct inspections and impose sanctions.
		\item Supervisors must use a risk-based approach to allocate resources and determine the frequency of on-site and off-site examinations.
		\item Five-year examination frequency is sufficient for all banks regardless of risk profile.
		\item Supervisors are not required to impose penalties for AML violations; only criminal courts have that authority.
		\item Off-site monitoring is optional under Recommendation 26.
	\end{enumerate}
	
	\textbf{Correct Answers: A and B}
	
	\subsection{19. FATF Recommendation 30: Law Enforcement Authorities \textnormal{[SELECT TWO]}}
	
	A country's police force has never conducted a parallel financial investigation in any drug trafficking case. All financial investigations are handled by a separate financial crimes unit that has no access to law enforcement intelligence.
	
	Which two statements correctly describe the country's obligations under Recommendation 30?
	
	\begin{enumerate}[label=\Alph*.]
		\item Law enforcement must have responsibility for conducting financial investigations into ML/TF.
		\item Law enforcement must be able to identify, trace, and initiate freezing and seizure of assets that are proceeds of crime.
		\item Financial investigations are optional and can be delegated to the FIU without law enforcement involvement.
		\item Asset freezing orders must be issued by a court; law enforcement cannot have provisional freezing authority.
		\item Only prosecutors have authority to conduct financial investigations; police have no role.
	\end{enumerate}
	
	\textbf{Correct Answers: A and B}
	
	\subsection{20. FATF Recommendation 32: Cash Courier \textnormal{[SELECT TWO]}}
	
	A country has no declaration system for cross-border physical transportation of currency. Customs agents have intercepted 14 cash couriers carrying over \$10 million in undeclared currency in the past year, but no penalties were imposed because the law does not require declarations.
	
	Which two statements correctly describe the country's obligations under Recommendation 32?
	
	\begin{enumerate}[label=\Alph*.]
		\item Countries must have a declaration or disclosure system for cross-border transportation of currency and bearer negotiable instruments.
		\item Countries must have the authority to stop or restrain cash couriers and impose sanctions for false declarations or non-disclosure.
		\item Cash courier controls are optional; Recommendation 32 is advisory only.
		\item Only cash amounts exceeding \$100,000 require declaration under FATF standards.
		\item Bearer negotiable instruments (traveler's checks, money orders) are exempt from declaration requirements.
	\end{enumerate}
	
	\textbf{Correct Answers: A and B}
	
	% ============================================================
	\section{FATF Mutual Evaluation Process and Immediate Outcomes}
	% ============================================================
	
	\subsection{21. FATF Mutual Evaluation: Two Pillars \textnormal{[SELECT TWO]}}
	
	A country is undergoing a FATF Mutual Evaluation. The assessment team notes strong technical compliance but poor effectiveness outcomes.
	
	Which two statements correctly describe how the country will be rated?
	
	\begin{enumerate}[label=\Alph*.]
		\item Countries are assessed on both Technical Compliance (TC) across all 40 Recommendations and Effectiveness across 11 Immediate Outcomes.
		\item A country can have high Technical Compliance but low Effectiveness, and this will be reflected in the overall assessment.
		\item Effectiveness is only assessed for Recommendations rated Non-Compliant; fully compliant Recommendations are presumed effective.
		\item Technical Compliance and Effectiveness ratings are combined into a single score; high TC compensates for low Effectiveness.
		\item The Mutual Evaluation process only assesses Technical Compliance; Effectiveness is assessed separately by the IMF.
	\end{enumerate}
	
	\textbf{Correct Answers: A and B}
	
	\subsection{22. FATF Immediate Outcomes (IO.1-IO.11) \textnormal{[SELECT THREE]}}
	
	A FATF assessor is reviewing the effectiveness of a country's AML/CFT system against the 11 Immediate Outcomes.
	
	Which three of the following correctly match an Immediate Outcome to its description?
	
	\begin{enumerate}[label=\Alph*.]
		\item \textbf{IO.1} — ML/TF risks are understood and, where appropriate, actions coordinated domestically.
		\item \textbf{IO.6} — Financial intelligence is used by competent authorities for ML/TF investigations.
		\item \textbf{IO.8} — Proceeds of crime and instrumentalities of crime are confiscated.
		\item \textbf{IO.3} — The FIU receives and disseminates STRs within 24 hours (this is not IO.3).
		\item \textbf{IO.11} — Financial institutions and DNFBPs are supervised for AML/CFT compliance.
	\end{enumerate}
	
	\textbf{Correct Answers: A, B, and C}
	
	\subsection{23. FATF Immediate Outcome 7 (IO.7): ML Investigations \textnormal{[SELECT ONE]}}
	
	A country has excellent STR filing rates (50,000 annually) and a well-staffed FIU. However, law enforcement has opened only three money laundering investigations in five years, and there have been zero convictions. The FIU's intelligence has never been used to initiate an investigation.
	
	What is the most likely Effectiveness rating for IO.7?
	
	\begin{enumerate}[label=\Alph*.]
		\item High — high STR volumes demonstrate effectiveness regardless of investigation outcomes.
		\item Moderate — some investigations were opened, so the system is functioning.
		\item Low — despite high STR volume, the lack of investigations and prosecutions indicates that financial intelligence is not being used to initiate or support ML investigations, which is the core of IO.7.
		\item Not applicable — IO.7 only applies to countries with a common law legal system.
		\item High — the FIU is receiving STRs, which is the only requirement for IO.7.
	\end{enumerate}
	
	\textbf{Correct Answer: C}
	
	\subsection{24. FATF Immediate Outcome 8 (IO.8): Confiscation \textnormal{[SELECT TWO]}}
	
	A country has laws authorizing criminal and civil forfeiture of proceeds of crime. However, in five years, no assets have been confiscated because prosecutors lack training on how to trace and identify criminal proceeds, and there is no asset management office to handle seized property.
	
	Which two findings would affect the IO.8 effectiveness rating?
	
	\begin{enumerate}[label=\Alph*.]
		\item The existence of confiscation laws alone does not satisfy IO.8; the country must demonstrate that proceeds are actually being identified, traced, frozen, and confiscated.
		\item The lack of prosecutor training on tracing and the absence of an asset management office are material deficiencies that will result in a Low effectiveness rating for IO.8.
		\item As long as the confiscation laws exist on paper, IO.8 is automatically rated High regardless of implementation.
		\item Confiscation only requires a conviction; asset tracing is not required for IO.8.
		\item Only criminal forfeiture counts for IO.8; civil forfeiture is not considered.
	\end{enumerate}
	
	\textbf{Correct Answers: A and B}
	
	% ============================================================
	\section{FATF Grey List, Black List, and Countermeasures}
	% ============================================================
	
	\subsection{25. FATF Grey List — Increased Monitoring \textnormal{[SELECT THREE]}}
	
	Jurisdiction Alpha has been placed on FATF's Increased Monitoring list (Grey List) due to strategic AML/CFT deficiencies. The country has committed to an action plan with FATF but has not yet fully addressed all deficiencies. A multinational bank maintains correspondent accounts with three banks in Jurisdiction Alpha.
	
	Which three statements correctly describe the bank's obligations and risk considerations?
	
	\begin{enumerate}[label=\Alph*.]
		\item The bank must apply Enhanced Due Diligence to all transactions and correspondent accounts involving Jurisdiction Alpha.
		\item Grey List designation requires EDD but does not automatically require termination of correspondent relationships.
		\item The bank should consider whether the EDD findings justify continuing, restricting, or terminating each relationship based on individual risk assessments.
		\item The bank must terminate all correspondent relationships with Grey List jurisdictions immediately.
		\item Grey List designation has no practical impact on financial institutions; it is purely a public statement with no required actions.
	\end{enumerate}
	
	\textbf{Correct Answers: A, B, and C}
	
	\subsection{26. FATF Black List — Call for Action \textnormal{[SELECT TWO]}}
	
	Jurisdiction Beta is on FATF's High-Risk Jurisdictions subject to a Call for Action (Black List). FATF has called on its members to apply countermeasures. A bank is evaluating its relationships with entities in Jurisdiction Beta.
	
	Which two statements correctly describe the bank's obligations?
	
	\begin{enumerate}[label=\Alph*.]
		\item FATF's call for countermeasures means the bank should evaluate and implement appropriate countermeasures, which may include severely restricting or terminating transactions and correspondent relationships.
		\item Countermeasures are mandatory for FATF members but are implemented through national laws and regulations; the bank must comply with its home country's implementation of FATF's call.
		\item Countermeasures are optional and can be ignored by private financial institutions.
		\item Banks must continue normal business relationships with Black List jurisdictions to promote financial inclusion.
		\item Black List designation has no legal effect; only UN Security Council sanctions require action.
	\end{enumerate}
	
	\textbf{Correct Answers: A and B}
	
	% ============================================================
	\section{USA PATRIOT Act Deep Dive}
	% ============================================================
	
	\subsection{27. USA PATRIOT Act Section 311: Primary ML Concern \textnormal{[SELECT THREE]}}
	
	The U.S. Secretary of the Treasury has designated a foreign financial institution as a "primary money laundering concern" under Section 311.
	
	Which three measures can be imposed on U.S. financial institutions regarding this foreign institution?
	
	\begin{enumerate}[label=\Alph*.]
		\item Prohibit U.S. financial institutions from opening or maintaining correspondent or payable-through accounts for the designated institution.
		\item Impose special recordkeeping and reporting requirements on transactions involving the designated institution.
		\item Impose conditions on correspondent accounts, such as requiring the foreign bank to pre-approve all transactions.
		\item Automatically freeze all assets of the foreign institution held in the United States without a court order.
		\item Seize all personal assets of the foreign institution's U.S.-based employees.
	\end{enumerate}
	
	\textbf{Correct Answers: A, B, and C}
	
	\subsection{28. USA PATRIOT Act Section 313: Shell Bank Prohibition \textnormal{[SELECT TWO]}}
	
	A U.S. correspondent bank discovers that a foreign respondent bank has no physical presence in any country, employs no staff, and is not affiliated with a regulated depository institution.
	
	Which two statements correctly describe the correspondent bank's obligations?
	
	\begin{enumerate}[label=\Alph*.]
		\item A "shell bank" is defined as a bank with no physical presence in any country and not affiliated with a regulated financial institution.
		\item U.S. banks are absolutely prohibited from maintaining correspondent accounts for foreign shell banks under Section 313.
		\item Shell banks are permitted if they are regulated by a G20 central bank.
		\item The correspondent bank may continue the relationship if it files a SAR every 90 days and increases monitoring.
		\item An exception exists for shell banks that are subsidiaries of U.S. banks.
	\end{enumerate}
	
	\textbf{Correct Answers: A and B}
	
	\subsection{29. USA PATRIOT Act Section 314(a): Law Enforcement Requests \textnormal{[SELECT TWO]}}
	
	A financial institution receives a FinCEN 314(a) request. The institution's compliance officer must implement the request.
	
	Which two statements correctly describe the institution's obligations under Section 314(a)?
	
	\begin{enumerate}[label=\Alph*.]
		\item The institution must search its accounts and transactions according to FinCEN's instructions for accounts maintained within the specified time period.
		\item The institution must maintain the confidentiality of the request and cannot disclose it to the subjects of the request.
		\item The institution must automatically freeze any matching account without further investigation.
		\item The institution must file a SAR for every account searched regardless of whether a match is found.
		\item The institution may ignore the request if the individuals are not current customers at the time of the request.
	\end{enumerate}
	
	\textbf{Correct Answers: A and B}
	
	\subsection{30. USA PATRIOT Act Section 314(b): Voluntary Information Sharing \textnormal{[SELECT TWO]}}
	
	Two U.S. banks have each identified suspicious activity involving different individuals who appear to be part of the same money laundering network. Neither bank alone has enough information to confirm the network's scope. Both banks have filed the required FinCEN opt-in notification for 314(b) sharing.
	
	Which two statements correctly describe their ability to share information?
	
	\begin{enumerate}[label=\Alph*.]
		\item Under Section 314(b), the banks may voluntarily share information about persons they suspect of ML/TF activity with each other.
		\item The sharing institutions receive a safe harbor from liability for the sharing.
		\item Section 314(b) requires both institutions to first obtain a court order before exchanging any customer information.
		\item Section 314(b) sharing automatically constitutes tipping-off because the shared information relates to a suspected subject.
		\item Section 314(b) sharing is only permissible if the amount of suspected ML exceeds \$1 million.
	\end{enumerate}
	
	\textbf{Correct Answers: A and B}
	
	\subsection{31. USA PATRIOT Act Section 319(b): Record Production \textnormal{[SELECT TWO]}}
	
	A U.S. bank maintains a correspondent account for a foreign bank. Federal law enforcement serves a subpoena on the U.S. bank demanding records related to the foreign bank's account. The foreign bank refuses to comply.
	
	Which two statements correctly describe Section 319(b)?
	
	\begin{enumerate}[label=\Alph*.]
		\item Section 319(b) allows the U.S. government to demand that a foreign bank produce records related to its U.S. correspondent account.
		\item If the foreign bank refuses to comply with a subpoena, the U.S. bank may be required to terminate the correspondent account, and funds in the account may be seized.
		\item Section 319(b) only applies to banks located in FATF Black List jurisdictions.
		\item Foreign banks are never required to produce records to U.S. authorities.
		\item Section 319(b) requires the U.S. bank to pay the foreign bank's legal fees for producing records.
	\end{enumerate}
	
	\textbf{Correct Answers: A and B}
	
	% ============================================================
	\section{EU Anti-Money Laundering Directives}
	% ============================================================
	
	\subsection{32. EU 4AMLD: Key Provisions \textnormal{[SELECT TWO]}}
	
	A European bank is reviewing its compliance with the Fourth Anti-Money Laundering Directive (4AMLD) as it prepares for implementation of 5AMLD.
	
	Which two provisions were introduced by 4AMLD?
	
	\begin{enumerate}[label=\Alph*.]
		\item Central registers of beneficial ownership of corporate entities (Beneficial Ownership Registers).
		\item Extension of CDD requirements to virtual currency exchange platforms and custodian wallet providers (this was 5AMLD, not 4AMLD).
		\item Requirement for member states to maintain central bank account registers (Bank Account Registers).
		\item Risk-based approach to AML/CFT supervision.
		\item Complete prohibition of anonymous prepaid cards.
	\end{enumerate}
	
	\textbf{Correct Answers: A and D}
	
	\subsection{33. EU 5AMLD: Key New Provisions \textnormal{[SELECT THREE]}}
	
	A European financial institution is updating its onboarding procedures to comply with 5AMLD.
	
	Which three provisions were introduced or expanded by 5AMLD?
	
	\begin{enumerate}[label=\Alph*.]
		\item Extension of CDD requirements to virtual currency exchange platforms and custodian wallet providers (VASPs).
		\item Publicly accessible central registries of beneficial owners for corporate entities (expanding 4AMLD).
		\item Extension of CDD requirements to prepaid card transactions above a reduced threshold (lowered from EUR 250 to EUR 150).
		\item Complete prohibition of all anonymous prepaid cards regardless of value.
		\item Creation of central bank account registers to allow FIUs to identify bank account holders.
	\end{enumerate}
	
	\textbf{Correct Answers: A, B, and C}
	
	\subsection{34. EU 6AMLD: Harmonized Predicate Offenses and Corporate Liability \textnormal{[SELECT TWO]}}
	
	A European bank's legal team is reviewing potential criminal exposure of a corporate client under 6AMLD.
	
	Which two provisions were introduced by 6AMLD?
	
	\begin{enumerate}[label=\Alph*.]
		\item Harmonization of 22 predicate offenses for money laundering across all EU member states.
		\item Extension of criminal liability to legal persons (corporations) for money laundering offenses.
		\item Elimination of all customer due diligence requirements for low-risk customers.
		\item Lowering of the wire transfer Travel Rule threshold from EUR 1,000 to EUR 500.
		\item Exemption of virtual assets from the scope of the Directive.
	\end{enumerate}
	
	\textbf{Correct Answers: A and B}
	
	% ============================================================
	\section{Wolfsberg Group, Basel Committee, and FSRBs}
	% ============================================================
	
	\subsection{35. Wolfsberg Group: Correspondent Banking Principles \textnormal{[SELECT TWO]}}
	
	A global correspondent bank is applying Wolfsberg Group principles when evaluating a new respondent bank.
	
	Which two factors are considered most critical in the Wolfsberg Correspondent Banking Due Diligence Questionnaire (CBDDQ)?
	
	\begin{enumerate}[label=\Alph*.]
		\item The respondent bank's AML/CFT policies, procedures, and internal controls.
		\item The respondent bank's customer base, particularly high-risk segments such as MSBs, NBFIs, and foreign PEPs.
		\item The respondent bank's CEO's personal investment portfolio.
		\item The respondent bank's office furniture and decoration preferences.
		\item The respondent bank's annual employee holiday party budget.
	\end{enumerate}
	
	\textbf{Correct Answers: A and B}
	
	\subsection{36. Basel Committee on Banking Supervision \textnormal{[SELECT TWO]}}
	
	The Basel Committee on Banking Supervision has issued guidance on the "Sound Management of Risks Related to ML and TF."
	
	Which two statements correctly describe the Basel Committee's expectations for AML/CFT risk management?
	
	\begin{enumerate}[label=\Alph*.]
		\item AML/CFT risk management should be integrated into the bank's overall operational risk management framework across all three lines of defense.
		\item Banks should conduct Enterprise-Wide Risk Assessments (EWRAs) to identify ML/TF risks across all products, services, geographies, and customer types.
		\item AML compliance should be isolated as a standalone silo function reporting only to the Chief Compliance Officer, with no integration into broader risk management.
		\item The first line of defense is internal audit; the third line of defense is business operations.
		\item The Basel Committee's guidance applies only to banks in FATF Grey List jurisdictions.
	\end{enumerate}
	
	\textbf{Correct Answers: A and B}
	
	\subsection{37. FATF-Style Regional Bodies (FSRBs) \textnormal{[SELECT TWO]}}
	
	A compliance officer is working with institutions in multiple regions and needs to understand the role of FSRBs.
	
	Which two statements correctly describe the function and authority of FSRBs?
	
	\begin{enumerate}[label=\Alph*.]
		\item FSRBs are autonomous regional bodies that promote, implement, and assess the execution of FATF Recommendations within their specific geographic regions.
		\item FSRBs conduct mutual evaluations of their member jurisdictions using the same methodology as FATF.
		\item FSRBs have no authority to assess member jurisdictions; only FATF can conduct mutual evaluations.
		\item FSRBs are subordinate operational branches staffed directly by FATF employees.
		\item FSRBs can impose sanctions on non-compliant jurisdictions independent of FATF.
	\end{enumerate}
	
	\textbf{Correct Answers: A and B}
	
	% ============================================================
	\section{UN Sanctions Framework}
	% ============================================================
	
	\subsection{38. UN Security Council Sanctions: Immediate Obligations \textnormal{[SELECT TWO]}}
	
	A bank processes a wire transfer from an entity that matches a name on the UN Security Council Consolidated Sanctions List. The compliance team must respond immediately.
	
	Which two actions are required?
	
	\begin{enumerate}[label=\Alph*.]
		\item Freeze the assets or funds immediately without prior notice to the entity, preventing any transfer, conversion, or movement of the asset.
		\item Report the freezing action to the relevant FIU and the UN Security Council Committee.
		\item Notify the customer of the freeze and give them 30 days to appeal.
		\item Process the transaction normally but file a SAR within 30 days.
		\item Return the funds to the originating bank without notifying any authority.
	\end{enumerate}
	
	\textbf{Correct Answers: A and B}
	
	\subsection{39. UNSCR 1540: Proliferation Financing \textnormal{[SELECT TWO]}}
	
	A trade finance bank is reviewing a transaction involving export of dual-use goods to a country under UNSC sanctions.
	
	Which two statements correctly describe the bank's obligations under UNSCR 1540?
	
	\begin{enumerate}[label=\Alph*.]
		\item Banks must establish effective controls to prevent financial support for proliferation of weapons of mass destruction.
		\item Banks must screen transactions for potential proliferation financing red flags, including exports of dual-use goods to sanctioned end-users.
		\item Proliferation financing controls only apply to banks in the exporting country; receiving country banks have no obligations.
		\item UNSCR 1540 only applies to governments, not to private financial institutions.
		\item Dual-use goods screening is always the responsibility of the exporter; banks have no independent obligation.
	\end{enumerate}
	
	\textbf{Correct Answers: A and B}
	
	% ============================================================
	\section{Egmont Group and FIU Cooperation}
	% ============================================================
	
	\subsection{40. Egmont Group: FIU-to-FIU Information Exchange \textnormal{[SELECT TWO]}}
	
	An investigator at Financial Intelligence Unit (FIU) Alpha has identified suspicious cross-border wire transfers that terminate in the jurisdiction of FIU Beta. The investigator needs additional information.
	
	Which two statements correctly describe FIU-to-FIU exchange under the Egmont Group framework?
	
	\begin{enumerate}[label=\Alph*.]
		\item FIUs can exchange financial intelligence directly using the Egmont Group's secure communication platform (Egmont Secure Web).
		\item Exchanged information must be used only for AML/CFT purposes and the receiving FIU must maintain confidentiality.
		\item Egmont Group information exchange can replace an MLAT for criminal prosecution purposes.
		\item Any FIU in the world can join the Egmont Group without any membership requirements.
		\item Information exchanged through Egmont can be used for any purpose, including commercial purposes.
	\end{enumerate}
	
	\textbf{Correct Answers: A and B}
	
	% ============================================================
	\section{Mutual Legal Assistance Treaties (MLATs)}
	% ============================================================
	
	\subsection{41. MLAT Process for Evidence Gathering \textnormal{[SELECT TWO]}}
	
	Country A has identified that \$30 million in bribery proceeds has been deposited into a bank account in Country B. Country A and Country B have an active Mutual Legal Assistance Treaty (MLAT). Country A's prosecutor needs bank records to support a criminal prosecution.
	
	Which two statements correctly describe the function of an MLAT?
	
	\begin{enumerate}[label=\Alph*.]
		\item MLATs provide a formal, legally binding mechanism for governments to cooperate in the exchange of evidence, judicial documents, and asset tracking/freezing across borders.
		\item An MLAT request is the proper channel for obtaining bank records that will be admissible in criminal proceedings in Country A's courts.
		\item MLATs allow Country A to send its police forces to arrest the bank manager in Country B without local authorization.
		\item Country A's FIU can obtain the same records through an Egmont Group request without an MLAT, and those records will be admissible in criminal court.
		\item MLATs are informal, non-binding agreements with no legal force.
	\end{enumerate}
	
	\textbf{Correct Answers: A and B}
	
	% ============================================================
	\section{Law Enforcement Cooperation and PPPs}
	% ============================================================
	
	\subsection{42. Public-Private Partnerships (PPPs) \textnormal{[SELECT TWO]}}
	
	A country has established a public-private partnership (PPP) to combat money laundering. The partnership includes representatives from the FIU, law enforcement, financial supervisors, and major banks.
	
	Which two statements correctly describe the purpose and function of AML/CFT PPPs?
	
	\begin{enumerate}[label=\Alph*.]
		\item PPPs improve intelligence sharing and detection of complex financial crime networks by enabling structured communication between government authorities and financial institutions.
		\item PPPs can help financial institutions better understand emerging typologies and help law enforcement prioritize financial intelligence.
		\item PPPs replace SAR filing obligations; banks in PPPs do not need to file SARs.
		\item PPPs eliminate the need for independent testing of AML programs.
		\item PPPs are only effective for cash-based ML; virtual assets cannot be addressed through PPPs.
	\end{enumerate}
	
	\textbf{Correct Answers: A and B}
	
	% ============================================================
	\section{Global Regulations — High-Difficulty Mixed}
	% ============================================================
	
	\subsection{43. Comparing FATF, Wolfsberg, and Basel \textnormal{[SELECT TWO]}}
	
	A compliance officer is explaining the different roles of FATF, Wolfsberg Group, and Basel Committee to a new analyst.
	
	Which two statements correctly distinguish these bodies?
	
	\begin{enumerate}[label=\Alph*.]
		\item FATF is an inter-governmental body that sets global AML/CFT standards (the 40 Recommendations) that countries commit to implement.
		\item The Wolfsberg Group is a private association of global banks that develops best practices and guidance, particularly for correspondent banking.
		\item FATF Recommendations are voluntary best practices with no enforcement mechanism.
		\item The Basel Committee sets legally binding sanctions that all countries must follow.
		\item Wolfsberg guidance is mandatory for all banks in FATF member countries.
	\end{enumerate}
	
	\textbf{Correct Answers: A and B}
	
	\subsection{44. FATF Recommendation 10: CDD Triggers \textnormal{[SELECT THREE]}}
	
	A compliance officer is training staff on when CDD must be conducted.
	
	Which three of the following are CDD trigger events under FATF Recommendation 10?
	
	\begin{enumerate}[label=\Alph*.]
		\item When establishing a business relationship.
		\item When carrying out an occasional transaction above the designated threshold (typically USD/EUR 15,000).
		\item When there is a suspicion of money laundering or terrorist financing, regardless of any threshold.
		\item When the customer's account balance drops below \$100.
		\item When the customer closes the account.
	\end{enumerate}
	
	\textbf{Correct Answers: A, B, and C}
	
	\subsection{45. FATF Recommendation 19: EDD for High-Risk Countries \textnormal{[SELECT TWO]}}
	
	A bank has identified a customer who is a national of a country on FATF's Grey List. The customer's business transactions involve wire transfers to that country.
	
	Which two statements correctly describe the bank's obligations?
	
	\begin{enumerate}[label=\Alph*.]
		\item The bank must apply Enhanced Due Diligence (EDD) to all transactions involving the high-risk country.
		\item EDD may include requiring additional information on the source of funds, the purpose of the transaction, and the beneficial owner.
		\item The bank must automatically refuse all transactions from Grey List countries.
		\item EDD is only required if the customer is a PEP; nationality alone is not a trigger.
		\item The bank must close all accounts held by nationals of Grey List countries.
	\end{enumerate}
	
	\textbf{Correct Answers: A and B}
	
	\subsection{46. FATF Recommendation 20: STR Filing Timing \textnormal{[SELECT ONE]}}
	
	A bank's AML analyst has identified a transaction that gives rise to reasonable suspicion of money laundering. The analyst has not yet completed a full investigation.
	
	Under Recommendation 20, when must the bank file the STR?
	
	\begin{enumerate}[label=\Alph*.]
		\item After completing a full investigation and obtaining conclusive proof of ML.
		\item Immediately upon suspicion, regardless of whether a full investigation is complete.
		\item Within 30 days of the suspicious transaction.
		\item Within 90 days of the suspicious transaction.
		\item Only after obtaining approval from the customer's relationship manager.
	\end{enumerate}
	
	\textbf{Correct Answer: B}
	
	\subsection{47. FATF Recommendation 24: Beneficial Ownership of Legal Persons \textnormal{[SELECT TWO]}}
	
	A country has enacted a law requiring all companies to identify their beneficial owners and file the information with a central registry. However, the registry is not accessible to financial institutions, law enforcement access requires a court order, and there is no requirement to update the information after initial filing. A FATF assessor notes these deficiencies.
	
	Which two statements correctly describe the deficiencies under Recommendation 24?
	
	\begin{enumerate}[label=\Alph*.]
		\item Beneficial ownership information must be accessible to competent authorities without a court order.
		\item There must be mechanisms to ensure that beneficial ownership information is accurate and updated.
		\item A central registry with no access for financial institutions satisfies Recommendation 24 as long as the information exists somewhere.
		\item Financial institutions must rely on customer declarations; access to government registries is optional.
		\item Beneficial ownership information only needs to be collected for companies with assets over \$1 million.
	\end{enumerate}
	
	\textbf{Correct Answers: A and B}
	
	\subsection{48. FATF Recommendation 25: Beneficial Ownership of Trusts \textnormal{[SELECT TWO]}}
	
	A country has many express trusts but does not require trustees to identify or maintain records of beneficial ownership. A FATF assessor identifies this as a deficiency.
	
	Which two statements correctly describe the obligations under Recommendation 25?
	
	\begin{enumerate}[label=\Alph*.]
		\item Trustees of express trusts must hold adequate, accurate, and current information on beneficial ownership of the trust.
		\item The information must be accessible to competent authorities.
		\item Trusts are completely exempt from all beneficial ownership requirements under FATF.
		\item Only trusts with assets exceeding EUR 1 million must comply with Recommendation 25.
		\item Trusts established in other countries are exempt from all requirements in the country where they operate.
	\end{enumerate}
	
	\textbf{Correct Answers: A and B}
	
	\subsection{49. FATF Recommendation 30: Financial Investigation Functions \textnormal{[SELECT TWO]}}
	
	A country's law enforcement has specialized financial investigation units within the police. However, these units do not have the power to obtain financial records directly from banks; all record requests must go through the prosecutor's office, which takes 6-12 months. The prosecutor's office has a backlog of 1,200 pending record requests.
	
	Which two statements correctly describe the deficiencies under Recommendation 30?
	
	\begin{enumerate}[label=\Alph*.]
		\item Law enforcement must have the power to obtain financial records expeditiously, without unreasonable delay.
		\item The 6-12 month delay in obtaining records undermines the effectiveness of financial investigations.
		\item A prosecutor-only approval system is acceptable as long as the law exists.
		\item Law enforcement does not need direct access to bank records; prosecutors are competent authorities.
		\item Financial investigations are optional; law enforcement can choose to investigate other crimes instead.
	\end{enumerate}
	
	\textbf{Correct Answers: A and B}
	
	\subsection{50. FATF Recommendation 32: Cross-Border Cash Courier Controls \textnormal{[SELECT TWO]}}
	
	A country has no declaration system for cross-border currency movement. Customs officers have intercepted 20 cash couriers carrying over \$15 million in undeclared currency in the past year, but the country imposes no penalties because the law does not require declarations. A FATF assessor identifies this as a deficiency.
	
	Which two statements correctly describe the obligations under Recommendation 32?
	
	\begin{enumerate}[label=\Alph*.]
		\item Countries must have a declaration or disclosure system for cross-border transportation of currency and bearer negotiable instruments.
		\item Countries must have the authority to stop or restrain cash couriers and impose sanctions for false declarations or non-disclosure.
		\item Cash courier controls are optional; Recommendation 32 is only a best practice.
		\item Only cash amounts exceeding \$100,000 require declaration under FATF standards.
		\item Bearer negotiable instruments (traveler's checks, money orders) are exempt from declaration requirements.
	\end{enumerate}
	
	\textbf{Correct Answers: A and B}
	
\end{document}